% Part: first-order-logic
% Chapter: syntax-and-semantics
% Section: terms-formulas

\documentclass[../../../include/open-logic-section]{subfiles}

\begin{document}

\olfileid{fol}{syn}{frm}

\olsection{ترمونه او \printtoken{P}{formula}}

يوه د لومړۍ درجې ژبه~$\Lang L$ چې راکړل شي، د~$\Lang L$ له بنسټيز
لغت څخه جوړ شوي عبارتونه تعريفولے شو۔ په دغو کښې په ځانګړي ډول
\emph{ترمونه} او \emph{!!{formula}s} شامل دي۔

\begin{defn}[ترمونه]
\ollabel{defn:terms}
د~$\Lang L$ د \emph{ترمونو} سټ~$\Trm[L]$ په استقرايي ډول داسې
تعريفېږي:
\begin{enumerate}
\item هر !!{variable} يو ترم دے۔
\item د~$\Lang L$ هر !!{constant} يو ترم دے۔
\item که $f$ يو $n$-ځايه !!{function} وي او $t_1$، \dots، $t_n$
  ترمونه وي، نو $\Atom{f}{t_1, \ldots, t_n}$ يو ترم دے۔
\tagitem{limitClause}{
  بل هېڅ څيز ترم نۀ دے۔}{}
\end{enumerate}
هغه ترم چې هېڅ !!{variable}s نۀ لري، \emph{تړلے ترم} بلل کېږي۔
\end{defn}

\begin{explain}
!!{constant}s زموږ د ژبې په مشخصاتو او د ترمونو په تعريف کښې د
سمبولونو د يوې جلا ډلې په توګه راځي، خو د صفر ځايه !!{function}s په
توګه هم ور شاملېدے شول۔ بيا د ترمونو د تعريف دويم شرط ته اړتيا نۀ
وه۔ يوازې دا بايد ومنو چې که $n = 0$ وي، نو
$\Atom{f}{t_1, \ldots, t_n}$ پخپله يوازې $f$ معنا لري۔
\end{explain}

\begin{defn}[فارمولونه]
\ollabel{defn:formulas}
د ژبې~$\Lang L$ د \emph{!!{formula}s} سټ~$\Frm[L]$ په استقرايي
ډول داسې تعريفېږي:
\begin{enumerate}
\tagitem{prvFalse}{$\lfalse$ يو اټومي !!{formula} دے۔}{}

\tagitem{prvTrue}{$\ltrue$ يو اټومي !!{formula} دے۔}{}

\item که $R$ د~$\Lang L$ يو $n$-ځايه !!{predicate} وي او $t_1$، \dots،
  $t_n$ د~$\Lang L$ ترمونه وي، نو $\Atom{R}{t_1,\ldots, t_n}$ يو
  اټومي !!{formula} دے۔

\item که $t_1$ او $t_2$ د~$\Lang L$ ترمونه وي، نو
  $\Atom{\eq}{t_1, t_2}$ يو اټومي !!{formula} دے۔

\tagitem{prvNot}{که $!A$ !!a{formula} وي، نو $\lnot !A$ هم
  !!a{formula} دے۔}{}

\tagitem{prvAnd}{که $!A$ او $!B$ !!{formula}s وي، نو $(!A \land
  !B)$ !!a{formula} دے۔}{}

\tagitem{prvOr}{که $!A$ او $!B$ !!{formula}s وي، نو $(!A \lor !B)$
  !!a{formula} دے۔}{}

\tagitem{prvIf}{که $!A$ او $!B$ !!{formula}s وي، نو $(!A \lif !B)$
  !!a{formula} دے۔}{}

\tagitem{prvIff}{که $!A$ او $!B$ !!{formula}s وي، نو $(!A \liff !B)$
  !!a{formula} دے۔}{}

\tagitem{prvAll}{که $!A$ !!a{formula} او $x$ !!a{variable} وي، نو
  $\lforall[x][!A]$ !!a{formula} دے۔}{}

\tagitem{prvEx}{که $!A$ !!a{formula} او $x$ !!a{variable} وي، نو
  $\lexists[x][!A]$ !!a{formula} دے۔}{}

\tagitem{limitClause}{بل هېڅ څيز !!a{formula} نۀ دے۔}{}
\end{enumerate}
\end{defn}

\begin{explain}
د ترمونو د سټ او د !!{formula}s د سټ تعريفونه \emph{استقرايي
تعريفونه} دي۔ په اصل کښې موږ د !!{formula}s سټ په نامتناهي ډېرو
پړاوونو کښې جوړوو۔ په لومړني پړاو کښې ټول اټومي فارمولونه فارمولونه
ګڼو؛ دا د تعريف له لومړنيو څو حالتونو سره سمون لري، يعنې د
\iftag{prvTrue}{$\ltrue$، }{}%
\iftag{prvFalse}{$\lfalse$، }{}%
$\Atom{R}{t_1,\dots,t_n}$ او $\Atom{\eq}{t_1,t_2}$ حالتونه۔ نو
``اټومي !!{formula}'' د همدې بڼو هر !!{formula} ته وايي۔

د تعريف نور حالتونه له هغو !!{formula}s څخه چې لا جوړ شوي وي، د نويو
!!{formula}s د جوړولو قاعدې راکوي۔ په دويم پړاو کښې د دغو قاعدو په
کارولو له اټومي !!{formula}s څخه !!{formula}s جوړولے شو۔ په درېيم
پړاو کښې له اټومي فارمولونو او په دويم پړاو کښې له ترلاسه شوو فارمولونو
څخه نوي فارمولونه جوړوو، او همداسې مخکښې۔ !!{formula} هر هغه څيز
دے چې بالاخره په داسې يوه پړاو کښې جوړ شي، او بل هېڅ څيز نۀ دے۔
\end{explain}

د دود له مخې، $\eq$ د خپلو مدخلو تر منځ ليکو او قوسونه پرېږدو:
$\eq[t_1][t_2]$ د $\Atom{\eq}{t_1,t_2}$ لنډون دے۔ همدارنګه
$\lnot \Atom{\eq}{t_1,t_2}$ د $\eq/[t_1][t_2]$ په بڼه لنډوو۔ کله
چې د دوه ځايه نښلوونکي~$\ast$ په کارولو له $!B$ او $!C$ څخه جوړ
شوے فارمول $(!B \ast !C)$ ليکو، ډېر ځله تر ټولو دباندې جوړه قوسونه
پرېږدو او يوازې~$!B \ast !C$ ليکو۔

\begin{intro}
د منطق ځينې متنونه د دې دپاره چې
\iftag{prvEx}{$\lexists[x][!A]$ }{}%
\iftag{notprvEx,notprvAll}{}{او }%
\iftag{prvAll}{$\lforall[x][!A]$ }{}%
\iftag{notprvEx,notprvAll}{!!a{formula}}{!!{formula}s}
وګڼل شي، دا شرط لګوي چې !!{variable}~$x$ بايد په~$!A$ کښې راشي۔
که دا شرط ونۀ لګوئ، کومه ستونزه نۀ پيدا کېږي او کار اسانه کېږي۔
\end{intro}

\begin{tagblock}{defNot,defOr,defAnd,defIf,defIff,defTrue,defFalse,defEx,defAll}
\begin{defn}
هغه فارمولونه چې په تعريف شوو عملګرونو جوړېږي، بايد په لاندې ډول
وپېژندل شي:

\begin{tagenumerate}{defTrue,defFalse,defNot,defOr,defAnd,defIf,defIff,defEx,defAll}

\tagitem{defTrue}{$\ltrue$ د
  \iftag{prvFalse}{$\lnot\lfalse$}{د کوم ثابت اټومي
    !!{formula}~$!A$ دپاره د $(!A \lor \lnot !A)$} لنډون دے۔}{}

\tagitem{defFalse}{$\lfalse$ د
  \iftag{prvTrue}{$\lnot\ltrue$}{د کوم ثابت اټومي
    !!{formula}~$!A$ دپاره د $(!A \land \lnot !A)$} لنډون دے۔}{}

\tagitem{defNot}{$\lnot !A$ د $!A \lif \lfalse$ لنډون دے۔}{}

\tagitem{defOr}{$!A \lor !B$ د
  \iftag{prvAnd}{$\lnot(\lnot !A \land \lnot !B)$}{$\lnot !A \lif
    !B$} لنډون دے۔}{}

\tagitem{defAnd}{$!A \land !B$ د
  \iftag{prvOr}{$\lnot(\lnot !A \lor \lnot !B)$}{$\lnot (!A \lif
    \lnot !B)$} لنډون دے۔}{}

\tagitem{defIf}{$!A \lif !B$ د
  \iftag{prvOr}{$\lnot !A \lor !B$}{$\lnot (!A \land \lnot !B)$}
  لنډون دے۔}{}
\textbf{د ژباړې د نسخې سمون (OLFOL-005):} سرچينه د لومړي لنډون
تر تړونکې رياضي نښې وروسته يو بې‌جوړه قوس راوړي؛ دلته زيات قوس لرې
شوے دے۔

\tagitem{defIff}{$!A \liff !B$ د $(!A \lif !B) \land (!B
  \lif !A)$ لنډون دے۔}{}

\tagitem{defAll}{$\lforall[x][!A]$ د $\lnot\lexists[x][\lnot !A]$
  لنډون دے۔}{}

\tagitem{defEx}{$\lexists[x][!A]$ د $\lnot\lforall[x][\lnot !A]$
  لنډون دے۔}{}
\end{tagenumerate}
\end{defn}
\end{tagblock}

کله چې د يوې ځانګړې کارونې په ژبه کښې کار کوو، ډېر ځله دوه ځايه
!!{predicate}s او !!{function}s د اړوندو ترمونو تر منځ ليکو؛ مثلاً د
حساب په ژبه کښې $t_1 < t_2$ او $(t_1 + t_2)$، او د سټ تيورۍ په ژبه
کښې $t_1 \in t_2$ ليکو۔ د حساب په ژبه کښې د ناستي تابع د دود له مخې
آن د خپل مدخل \emph{وروسته} ليکل کېږي:~$t'$۔ خو په رسمي ډول دا په
ترتيب سره يوازې د $\Atom{\Obj A^2_0}{t_1, t_2}$،
$\Obj f^2_0(t_1, t_2)$، $\Atom{\Obj A^2_0}{t_1, t_2}$ او
$\Obj f^1_0(t)$ دوديز لنډونونه دي۔
\textbf{د ژباړې د نسخې سمون (OLFOL-006):} سرچينه د ناستي تابع په
رسمي سمبول کښې د څيز نښه پرېږدي؛ دلته هغه د همدې فهرست له سمبوليزم
سره سمه ور زياته شوې ده۔

\begin{defn}[نحوي يوشانوالے]
سمبول $\ident$ د سمبولونو د تارونو تر منځ نحوي يوشانوالے څرګندوي؛
يعنې $!A \ident !B$ هله او يوازې هله چې $!A$ او $!B$ د يو شان
اوږدوالي لرونکي د سمبولونو تارونه وي او په هر ځاے کښې يو شان سمبول
ولري۔
\end{defn}

د $\ident$ سمبول دواړو خواوو ته هغه تارونه هم ليکل کېدے شي چې په نښلولو
ترلاسه شوي وي؛ مثلاً $!A \ident (!B \lor !C)$ دا معنا لري: د
سمبولونو تار~$!A$ له هغه تار سره يو شان دے چې په همدې ترتيب د پرانيستي
قوس، تار $!B$، سمبول $\lor$، تار~$!C$ او تړلي قوس له نښلولو ترلاسه
کېږي۔ که دا حالت وي، نو پوهېږو چې د~$!A$ لومړے سمبول پرانيستے قوس
دے، $!A$ د دويم سمبول له ځايه $!B$ د يوه فرعي تار په توګه لري، له
هغه فرعي تار وروسته~$\lor$ راځي، او همداسې نور۔

څرنګه چې ترمونه او !!{formula}s د استقرايي تعريفونو له لارې له بنسټيزو
اجزاوو څخه جوړېږي، د هغو په باب د څه ثابتولو دپاره لاندې استقرايي
اصول کارولے شو۔

\begin{lem}[\emph{پر ترمونو د استقرا اصل}]
    \ollabel{lem:trmind}
    فرض کړئ $\Lang L$ د لومړۍ درجې يوه ژبه ده۔
    که يو خاصيت~$P$ داسې وي چې
    %
    \begin{enumerate}
        \item د هر !!{variable}~$v$ دپاره صدق کوي؛
        %
        \item د~$\Lang L$ د هر !!{constant}~$a$ دپاره صدق کوي؛ او
        %
        \item کله چې د $t_1$،~\dots، $t_n$ دپاره صدق کوي او $f$
        د~$\Lang L$ يو $n$-ځايه !!{function} وي، د
        $f(t_1,\dotsc,t_n)$ دپاره هم صدق کوي،
    \end{enumerate}
    (په دې فرض چې $t_1$،~\dots، $t_n$ د~$\Lang{L}$ ترمونه دي)،
    نو $P$ د~$\Trm[L]$ د هر ترم دپاره صدق کوي۔
\end{lem}

\begin{prob}
    \olref[fol][syn][frm]{lem:trmind} ثابته کړئ۔
\end{prob}

\begin{lem}[\emph{پر !!{formula}s د استقرا اصل}]
    \ollabel{thm:frmind}
    فرض کړئ $\Lang L$ د لومړۍ درجې يوه ژبه ده۔
    که يو خاصيت~$P$ د ټولو اټومي !!{formula}s دپاره صدق کوي او داسې
    وي چې
    %
    \begin{enumerate}
        \tagitem{prvNot}{کله چې د~$!A$ دپاره صدق کوي، د
            $\lnot !A$ دپاره هم صدق کوي؛}{}
        \tagitem{prvAnd}{کله چې د $!A$ او~$!B$ دپاره صدق کوي، د
            $(!A \land !B)$ دپاره هم صدق کوي؛}{}
        \tagitem{prvOr}{کله چې د $!A$ او~$!B$ دپاره صدق کوي، د
            $(!A \lor !B)$ دپاره هم صدق کوي؛}{}
        \tagitem{prvIf}{کله چې د $!A$ او~$!B$ دپاره صدق کوي، د
            $(!A \lif !B)$ دپاره هم صدق کوي؛}{}
        \tagitem{prvIff}{کله چې د $!A$ او~$!B$ دپاره صدق کوي، د
            $(!A \liff !B)$ دپاره هم صدق کوي؛}{}
        \tagitem{prvEx}{کله چې د~$!A$ دپاره صدق کوي، د
            $\lexists[x][!A]$ دپاره هم صدق کوي؛}{}
        \tagitem{prvAll}{کله چې د~$!A$ دپاره صدق کوي، د
            $\lforall[x][!A]$ دپاره هم صدق کوي؛}{}
  \end{enumerate}
  (په دې فرض چې $!A$ او $!B$ د~$\Lang{L}$ !!{formula}s دي)،
  نو $P$ د~$\Frm[L]$ د ټولو فارمولونو دپاره صدق کوي۔
\end{lem}

\end{document}
